\input{preamble}
\usepackage{amssymb}

\newcommand{\rightlooparrow}{\mathbin{
\vbox{\openup-10.25pt\halign{\hss$##$\hss\cr
\circ\cr\longrightarrow\cr}}
}}

\begin{document}

\title{Differential Topology}
\maketitle

\phantomsection
\label{section-phantom}

\tableofcontents

\section{Transversality theorems}
\label{section-transversality-theorems}

\noindent
A simple case of the following definition is
when $f$ is the inclusion and $X,W$ are both
submanifolds of $Y$.

\begin{definition}
\label{definition-transversal-map}
Let $X,Y$ be smooth manifolds and let
$f \in C^\infty(X,Y)$.
Let $W \subset Y$ be a submanifold.
We say that $f$ is {\it transverse} to $W$,
and write $f \pitchfork W$, if for every
$x \in f^{-1}(W)$ we have
$$
d_xf(T_xX)+T_{f(x)}W=T_{f(x)}Y.
$$
\end{definition}

\noindent
This imposes restrictions on the dimensions
of $X$ and $W$:
$$
\dim X+\dim W \geq \dim Y.
$$
This may also be thought of as a
generalization of the regular value
condition.  If $W$ is a single point $y$,
then $f \pitchfork \{y\}$ says exactly that
$$
d_xf(T_xX)=T_yY,
\qquad x \in f^{-1}(y),
$$
that is, $df$ is surjective at every point of
the fiber.

\begin{theorem}
\label{theorem-preimage-transverse}
Let $f:X \to Y$ be smooth and let
$W \subset Y$ be a submanifold.  If
$f \pitchfork W$, then $f^{-1}(W)$ is a
submanifold of $X$ and
$$
\operatorname{codim}_X f^{-1}(W)
=
\operatorname{codim}_Y W.
$$
\end{theorem}

\begin{proof}
This is the regular value theorem in local
coordinates.  Near a point of $W$ choose
coordinates in which $W$ is cut out by the
last $c$ coordinates, where
$c=\operatorname{codim}W$.  The map to these
last $c$ coordinates has surjective
differential along $f^{-1}(W)$ exactly because
of transversality.  Hence its zero set is a
submanifold of codimension $c$.
\end{proof}

\section{Course plan and review}
\label{section-course-plan-review}

\begin{reference}
These notes follow Vinicius Ramos' course on
differential topology.  Standard references
are \cite{milnordt}, \cite{muk}, and
\cite{hirsch}.
\end{reference}

\noindent
The course plan was the following.

\begin{enumerate}
\item Review of manifolds.

\item Transversality: Sard's theorem,
Whitney topologies, and approximation.

\item Intersection theory and index.

\item Morse theory.

\item Additional possible topics:
$h$-cobordism, low-dimensional topology, and
Poincar\'e in dimensions $n\geq 5$.
\end{enumerate}

\noindent
We use the standard convention that a smooth
manifold is Hausdorff, second countable, and
locally Euclidean, equipped with a maximal
smooth atlas.  The dimension is locally
constant, hence constant on connected
components.  There are topological atlases
which do not contain a smooth subatlas.

\begin{theorem}
\label{theorem-smoothing-manifolds}
Every $C^k$-manifold, for
$k=1,\ldots,\infty$, is $C^k$-diffeomorphic
to a $C^\infty$-manifold.
\end{theorem}

\begin{proof}
We record the statement as one of the basic
smoothing theorems used in the course.  The
proof is not included here.
\end{proof}

\begin{theorem}
\label{theorem-compatible-smoothings}
Let $1\leq \ell \leq k \leq \infty$.
If two $C^k$-manifolds are
$C^\ell$-diffeomorphic, then they are
$C^k$-diffeomorphic.
\end{theorem}

\begin{proof}
We omit the proof.  The point is that lower
regularity compatibility is enough to compare
the induced smooth structures.
\end{proof}

\begin{theorem}
\label{theorem-partition-of-unity}
Let $M$ be a smooth manifold and let
$\{U_i\}$ be an open cover.  Then there is a
smooth partition of unity subordinate to this
cover.
\end{theorem}

\begin{proof}
This uses paracompactness of smooth
manifolds.  Every topological manifold is
paracompact, and smooth bump functions allow
one to refine locally finite covers by smooth
functions.
\end{proof}

\section{Tangent and cotangent bundles}
\label{section-tangent-cotangent-bundles}

\noindent
Let $M$ be a smooth manifold.  The tangent
space $T_pM$ can be defined either as
velocities of curves through $p$ or as
derivations at $p$.  In a chart
$(U,\varphi)$ with
$\varphi=(x^1,\ldots,x^n)$, a basis is
$$
\left\{
\frac{\partial}{\partial x^1}\Big|_p,
\ldots,
\frac{\partial}{\partial x^n}\Big|_p
\right\}.
$$
The cotangent space is
$$
T_p^*M=(T_pM)^*
=\operatorname{Hom}(T_pM,\mathbb R),
$$
with dual basis
$$
\left\{dx^1|_p,\ldots,dx^n|_p\right\},
\qquad
 dx^i|_p
\left(
\frac{\partial}{\partial x^j}\Big|_p
\right)=\delta^i_j.
$$
Changing charts changes these bases; there is
no canonical basis of $T_pM$ or $T_p^*M$.

\begin{exercise}
\label{exercise-change-of-basis}
Let $(U,\varphi)$ and $(V,\psi)$ be charts
near $p$, with coordinates
$x^1,\ldots,x^n$ and $y^1,\ldots,y^n$.
Show that
$$
\frac{\partial}{\partial x^j}
=
\sum_{i=1}^n
\frac{\partial y^i}{\partial x^j}
\frac{\partial}{\partial y^i}.
$$
\end{exercise}

\noindent
The tangent bundle is
$$
TM=\bigsqcup_{p\in M}T_pM.
$$
For every chart $(U,\varphi)$ there is a
bijection
$$
\begin{aligned}
U\times \mathbb R^n
&\longrightarrow \pi^{-1}(U)\\
(p,(v_1,\ldots,v_n))
&\longmapsto
\sum_{i=1}^n v_i
\frac{\partial}{\partial x^i}\Big|_p.
\end{aligned}
$$
Using these bijections one topologizes $TM$
and gives it a smooth structure of dimension
$2n$.  The same construction gives the
cotangent bundle
$$
T^*M=\bigsqcup_{p\in M}T_p^*M.
$$
If $M$ is only $C^k$, then $TM$ is naturally
$C^{k-1}$.

\begin{definition}
\label{definition-parallelizable}
A smooth $n$-manifold $M$ is
{\it parallelizable} if
$$
TM\cong M\times \mathbb R^n.
$$
\end{definition}

\begin{example}
\label{example-parallelizable-manifolds}
The manifolds $\mathbb R^n$ and $S^1$ are
parallelizable.  Every orientable
3-manifold is also parallelizable, although
that result is much deeper.
\end{example}

\section{Immersions, submersions, and embeddings}
\label{section-immersions-submersions-embeddings}

\begin{definition}
\label{definition-differential-map}
Let $f:M\to N$ be smooth.  The
{\it derivative} or {\it differential} of
$f$ at $p$ is the linear map
$$
Df_p:T_pM\to T_{f(p)}N.
$$
In the curve model of tangent vectors,
$Df_p[\gamma]=[f\circ\gamma]$.  In the
derivation model,
$$
(Df_pv)(g)=v(g\circ f),
\qquad g\in C^\infty(N).
$$
\end{definition}

\noindent
In local coordinates $x^j$ on $M$ and $y^i$
on $N$, this is represented by the Jacobian
matrix:
$$
Df_p
\left(\frac{\partial}{\partial x^j}\Big|_p
\right)
=
\sum_i
\frac{\partial f_i}{\partial x^j}(p)
\frac{\partial}{\partial y^i}\Big|_{f(p)}.
$$

\begin{definition}
\label{definition-immersion-submersion-embedding}
Let $f:M\to N$ be smooth.
\begin{enumerate}
\item $f$ is an {\it immersion} at $p$ if
$Df_p$ is injective.

\item $f$ is a {\it submersion} at $p$ if
$Df_p$ is surjective.

\item $f$ is an {\it embedding} if it is an
injective immersion and the inverse
$f(M)\to M$ is continuous, where $f(M)$ has
the subspace topology from $N$.
\end{enumerate}
\end{definition}

\noindent
We often write
$$
M\rightlooparrow N,
\qquad
M\hookrightarrow N,
\qquad
M\twoheadrightarrow N
$$
for immersions, embeddings, and submersions.
The continuity condition in the definition of
embedding is necessary: an injective immersed
curve may accumulate on itself in the target,
so that the inverse from its image is not
continuous.

\begin{remark}
\label{remark-embedding-submanifold}
If $f:M\to N$ is an embedding, then $f(M)$
inherits a unique smooth structure such that
$f:M\to f(M)$ is a diffeomorphism.  Thus
embeddings are precisely the maps whose
images are honest submanifolds.
\end{remark}

\begin{exercise}
\label{exercise-same-embedded-image}
Let $N_1,N_2$ be manifolds and let
$\varphi_i:N_i\to M$ be smooth embeddings.
Suppose their images coincide.  Show that
$N_1$ and $N_2$ are diffeomorphic.
\end{exercise}

\section{Regular values}
\label{section-regular-values}

\begin{definition}
\label{definition-regular-value}
Let $f:M\to N$ be smooth.  A point $y\in N$
is a {\it regular value} of $f$ if $f$ is a
submersion at every point of $f^{-1}(y)$.
Equivalently, $Df_x$ is surjective for all
$x\in f^{-1}(y)$.
\end{definition}

\begin{theorem}[Regular value theorem]
\label{theorem-regular-value}
If $y$ is a regular value of $f:M\to N$, then
$f^{-1}(y)$ is a submanifold of $M$ of
dimension
$$
\dim M-\dim N,
$$
provided the fiber is nonempty.
\end{theorem}

\begin{proof}
This is the implicit function theorem.  Take
charts around $x\in f^{-1}(y)$ and around
$y$, both centered at the relevant points.
In coordinates we get a map
$\Phi:\mathbb R^m\to\mathbb R^n$ with
surjective derivative at $0$.  After a local
change of coordinates, $\Phi$ becomes the
projection
$$
\mathbb R^n\times \mathbb R^{m-n}
\longrightarrow \mathbb R^n.
$$
Hence $\Phi^{-1}(0)$ is locally
$\{0\}\times \mathbb R^{m-n}$.
\end{proof}

\section{Vector bundles and sections}
\label{section-vector-bundles-sections}

\begin{definition}
\label{definition-vector-bundle}
Let $E$ and $M$ be manifolds and let
$\pi:E\to M$ be a surjective submersion.
We say that $\pi$ is a {\it vector bundle} of
rank $r$ if every fiber
$E_p=\pi^{-1}(p)$ is a vector space and for
every $p\in M$ there is a neighborhood $U$
and a diffeomorphism
$$
\varphi:\pi^{-1}(U)\to U\times\mathbb R^r
$$
which is linear on each fiber and commutes
with the projections to $U$.
\end{definition}

\begin{example}
\label{example-basic-vector-bundles}
The tangent bundle, cotangent bundle, and the
trivial bundle $M\times\mathbb R^r\to M$ are
vector bundles.
\end{example}

\begin{definition}
\label{definition-section-vector-bundle}
A {\it section} of a vector bundle
$\pi:E\to M$ is a smooth map
$s:M\to E$ such that $\pi\circ s=\text{id}_M$.
\end{definition}

\noindent
A vector field is a section of $TM$.  A
1-form is a section of $T^*M$.  If a section
is compactly supported, its support has
compact closure.  Partitions of unity are the
basic tool for producing global sections from
local data.

\section{Sard theorem}
\label{section-sard-theorem}

\begin{definition}
\label{definition-critical-point}
Let $f:X^m\to Y^n$ be smooth.  A point
$x\in X$ is a {\it critical point} if
$Df_x$ does not have maximal rank.  A point
$y\in Y$ is a {\it critical value} if
$y=f(x)$ for some critical point $x$.
\end{definition}

\begin{theorem}[Sard]
\label{theorem-sard}
Let $f:X\to Y$ be smooth.  The set of
critical values of $f$ has measure zero in
$Y$.  In particular, the regular values of
$f$ are dense.
\end{theorem}

\begin{proof}
We only record the structure of the proof.
In Euclidean coordinates one considers the
sets where all partial derivatives up to a
certain order vanish.  Taylor's theorem shows
that the image of each such set can be
covered by boxes whose total volume is
arbitrarily small.  The manifold case follows
by covering the domain with countably many
coordinate charts.
\end{proof}

\begin{remark}
\label{remark-sard-finite-regularity}
Sard's theorem also holds for $C^\ell$ maps
provided
$$
\ell>\max(m-n,0).
$$
This regularity threshold is important in the
full statement.
\end{remark}

\section{Jets}
\label{section-jets}

\begin{definition}
\label{definition-jet-equivalence}
Let $f,g:X\to Y$ be smooth maps.  We say
that $f$ and $g$ have the same $k$-jet at
$x\in X$ if, in local coordinates around $x$
and around $f(x)=g(x)$, all partial
derivatives up to order $k$ agree at $x$.
The equivalence class is denoted $j_x^kf$.
\end{definition}

\begin{definition}
\label{definition-jet-space}
The {\it $k$-jet space} $J^k(X,Y)$ is the set
of all $k$-jets $j_x^kf$ of smooth maps
$X\to Y$.
There are natural source and target maps
$$
\alpha:J^k(X,Y)\to X,
\qquad
\beta:J^k(X,Y)\to Y.
$$
\end{definition}

\noindent
The space $J^k(X,Y)$ is a smooth manifold.
For $k=0$ one has $J^0(X,Y)=X\times Y$.
For $k=1$, a point of $J^1(X,Y)$ consists of
$x\in X$, $y\in Y$, and a linear map
$T_xX\to T_yY$.
Thus
$$
J^1(X,Y)
\cong
\operatorname{Hom}(TX,TY)
$$
in the appropriate fibered sense.

\begin{definition}
\label{definition-jet-extension}
For a smooth map $f:X\to Y$, its
{\it $k$-jet extension} is the map
$$
j^kf:X\to J^k(X,Y),
\qquad
x\longmapsto j_x^kf.
$$
\end{definition}

\section{Whitney topologies}
\label{section-whitney-topologies}

\noindent
The space $C^\infty(X,Y)$ admits several
natural topologies.  The weak Whitney
topology controls a map and finitely many
of its derivatives on compact subsets.  The
strong Whitney topology allows the control to
vary over locally finite families of compact
sets.  These topologies are designed so that
statements such as ``transverse maps are
dense'' become precise.

\begin{definition}
\label{definition-delta-ball}
Let $\delta:X\to (0,\infty)$ and let
$f\in C^\infty(X,Y)$, after choosing a
metric on a jet bundle.  A $\delta$-ball
around $f$ consists of those maps $g$ whose
relevant jets stay within distance
$\delta(x)$ from the jets of $f$ at every
$x$.
\end{definition}

\begin{definition}
\label{definition-baire-space}
A topological space is a {\it Baire space} if
a countable intersection of open dense sets
is dense.
\end{definition}

\noindent
The strong Whitney topology is a Baire
topology on $C^\infty(X,Y)$ in the usual
situations.  This is why genericity
statements can be proved by intersecting
countably many open dense conditions.

\section{Parametric transversality}
\label{section-parametric-transversality}

\begin{theorem}[Parametric transversality]
\label{theorem-parametric-transversality}
Let $F:X\times B\to Y$ be smooth and let
$W\subset Y$ be a submanifold.  If
$F\pitchfork W$, then for almost every
$b\in B$ the map
$$
F_b:X\to Y,
\qquad
F_b(x)=F(x,b),
$$
is transverse to $W$.
\end{theorem}

\begin{proof}
Let $Q=F^{-1}(W)$.  Since $F\pitchfork W$,
Theorem \ref{theorem-preimage-transverse}
shows that $Q$ is a submanifold of
$X\times B$.  Consider the projection
$\pi:Q\to B$.  By Sard's theorem, almost
every $b\in B$ is a regular value of $\pi$.
For such $b$, the condition that $b$ be a
regular value of $\pi$ is equivalent to
$F_b\pitchfork W$.
\end{proof}

\begin{theorem}[Thom transversality]
\label{theorem-thom-transversality}
Let $X,Y$ be smooth manifolds and let
$W\subset J^k(X,Y)$ be a submanifold.  Then
$$
\{f\in C^\infty(X,Y):j^kf\pitchfork W\}
$$
is residual in $C^\infty(X,Y)$, and in the
strong Whitney topology it is dense.  If
$W$ is closed, this set is also open.
\end{theorem}

\begin{proof}
The proof reduces locally to the parametric
transversality theorem.  One perturbs a map
by adding polynomial terms in coordinate
balls.  The finite-dimensional space of
polynomial perturbations supplies the
parameter space $B$, and the evaluation map
from $X\times B$ to the jet space is made
transverse to $W$.  A partition of unity and
a Baire category argument globalize the
construction.
\end{proof}

\section{Multijets and Whitney theorems}
\label{section-multijets-whitney-theorems}

\noindent
Multijet spaces keep track of jets of a map
at several distinct points at once.  They are
used to impose conditions on pairs or tuples
of points, for example to avoid double points
or to control self-intersections.

\begin{theorem}[Whitney immersion theorem]
\label{theorem-whitney-immersion}
Every smooth $n$-manifold admits an immersion
into $\mathbb R^{2n-1}$.
\end{theorem}

\begin{proof}
Start with an embedding into some
$\mathbb R^N$.  Project generically to a
lower-dimensional Euclidean space.  The bad
projections are those for which the
differential drops rank somewhere; this is a
transversality condition on the 1-jet.  Thom
transversality shows that a generic
projection avoids the bad locus as long as
the target has dimension at least $2n-1$.
\end{proof}

\begin{theorem}[Whitney embedding theorem]
\label{theorem-whitney-embedding}
Every smooth $n$-manifold admits an embedding
into $\mathbb R^{2n}$.
\end{theorem}

\begin{proof}
Again start from an embedding into
$\mathbb R^N$ and project generically.  The
immersion condition is controlled by the
1-jet.  Injectivity is controlled by the
2-point multijet: one must avoid the diagonal
in $\mathbb R^{2n}\times\mathbb R^{2n}$ for
pairs of distinct points.  The relevant bad
sets have sufficiently large codimension for
generic projections to avoid them.
\end{proof}

\section{Manifolds with boundary}
\label{section-manifolds-with-boundary}

\begin{definition}
\label{definition-half-space}
Let
$$
\mathbb H^n=\{(x_1,\ldots,x_n)\in
\mathbb R^n:x_n\geq 0\}.
$$
A {\it smooth manifold with boundary} is a
space locally modeled on open subsets of
$\mathbb H^n$ with smooth transition maps.
\end{definition}

\noindent
The boundary $\partial M$ consists of those
points which in a half-space chart map to
$x_n=0$.  The interior consists of the points
which map to $x_n>0$.  This is independent
of the chart.

\begin{theorem}[Sard with boundary]
\label{theorem-sard-boundary}
Let $f:X\to Y$ be smooth, where $X$ is a
manifold with boundary.  If $a\in Y$ is a
regular value of both $f$ and
$f|_{\partial X}$, then $f^{-1}(a)$ is a
manifold with boundary and
$$
\partial(f^{-1}(a))=f^{-1}(a)\cap \partial X.
$$
\end{theorem}

\begin{proof}
Apply the regular value theorem in boundary
charts.  In interior charts this is the usual
argument.  In boundary charts the local model
is the intersection of a regular level set
with a half-space.
\end{proof}

\begin{theorem}[Thom transversality with boundary]
\label{theorem-thom-boundary}
Let $X$ be a manifold with boundary and let
$W\subset Y$ be a submanifold.  Generically,
a smooth map $f:X\to Y$ is transverse to
$W$ and $f|_{\partial X}$ is also transverse
to $W$.
\end{theorem}

\begin{proof}
Apply the ordinary Thom transversality
theorem simultaneously to $X$ and to
$\partial X$.  The boundary condition is an
extra transversality condition on the
restriction.
\end{proof}

\section{Orientation}
\label{section-orientation}

\begin{definition}
\label{definition-orientation-vector-space}
An {\it orientation} of a finite-dimensional
real vector space is a choice of one of the
two connected components of the set of
ordered bases.  Equivalently, two ordered
bases define the same orientation if the
change-of-basis determinant is positive.
\end{definition}

\begin{definition}
\label{definition-orientation-manifold}
An {\it orientation} of a smooth manifold
$M$ is a smoothly varying choice of
orientation on each tangent space $T_pM$.
\end{definition}

\noindent
If $M$ is an oriented manifold with boundary,
then $\partial M$ receives the boundary
orientation by the ``outward normal first''
rule: an ordered basis $(v_1,\ldots,v_{n-1})$
of $T_p\partial M$ is positive if
$(\nu,v_1,\ldots,v_{n-1})$ is a positive
basis of $T_pM$, where $\nu$ is an outward
normal vector.

\section{Intersection theory}
\label{section-intersection-theory}

\begin{definition}
\label{definition-intersection-number}
Let $X^n$ and $Y^m$ be oriented manifolds,
let $W^k\subset Y$ be an oriented
submanifold, and let $f:X\to Y$ be smooth and
transverse to $W$.  If
$$
\dim X+\dim W=\dim Y,
$$
then $f^{-1}(W)$ is a finite set when $X$ is
compact.  The {\it intersection number}
$I(f,W)$ is the signed count of this set.
\end{definition}

\noindent
The sign at a point $x\in f^{-1}(W)$ is
determined by comparing the orientation of
$T_{f(x)}Y$ with the orientation obtained
from
$$
df_x(T_xX)\oplus T_{f(x)}W.
$$

\begin{proposition}
\label{proposition-intersection-homotopy}
If $f_0,f_1:X\to Y$ are homotopic through
maps transverse to $W$ at the ends, then
$$
I(f_0,W)=I(f_1,W).
$$
\end{proposition}

\begin{proof}
Let $F:X\times[0,1]\to Y$ be a homotopy
between $f_0$ and $f_1$, perturbed relative
to the boundary so that $F\pitchfork W$.
Then $F^{-1}(W)$ is an oriented compact
1-manifold whose boundary is
$f_1^{-1}(W)-f_0^{-1}(W)$.  The signed number
of boundary points of a compact oriented
1-manifold is zero.
\end{proof}

\section{Degree}
\label{section-degree}

\begin{definition}
\label{definition-degree-map}
Let $X$ and $Y$ be compact oriented
$n$-manifolds, with $Y$ connected.  For a
smooth map $f:X\to Y$ and a regular value
$y\in Y$, the {\it degree} of $f$ is
$$
\deg(f)=\sum_{x\in f^{-1}(y)}\operatorname{sign}
(Df_x).
$$
\end{definition}

\begin{proposition}
\label{proposition-degree-well-defined}
The degree is independent of the chosen
regular value and invariant under homotopy.
\end{proposition}

\begin{proof}
This is the intersection number of $f$ with
the point $y\subset Y$.  Independence of $y$
and homotopy invariance follow from the
homotopy invariance of intersection number.
\end{proof}

\begin{example}
\label{example-degree-circle}
For maps $S^1\to S^1$, the degree is the
winding number.  If $f:S^1\to S^1$ is given
by $f(z)=z^d$, then $\deg(f)=d$.
\end{example}

\section{Fundamental theorem of algebra}
\label{section-fundamental-theorem-algebra-degree}

\begin{proposition}
\label{proposition-fta-degree}
Every nonconstant complex polynomial has a
complex root.
\end{proposition}

\begin{proof}
Let $p(z)$ be a polynomial of degree $d>0$.
For $R$ large, the map
$$
S_R^1\to S^1,
\qquad
z\longmapsto \frac{p(z)}{|p(z)|}
$$
has the same degree as
$z\mapsto z^d/|z|^d$, namely $d$.  If $p$ had
no root in the disk, this boundary map would
extend to the disk and hence would have
degree zero.  This is impossible.
\end{proof}

\section{Winding number and Jordan--Brouwer}
\label{section-winding-jordan-brouwer}

\noindent
For a smooth closed curve
$\gamma:S^1\to \mathbb R^2\setminus\{p\}$,
the winding number about $p$ is the degree of
$$
S^1\to S^1,
\qquad
z\longmapsto
\frac{\gamma(z)-p}{\|\gamma(z)-p\|}.
$$
It is constant on connected components of the
complement of the curve.

\begin{theorem}[Jordan--Brouwer separation]
\label{theorem-jordan-brouwer}
If $M\subset S^n$ is a compact connected
embedded hypersurface, then
$S^n\setminus M$ has exactly two connected
components, and $M$ is the boundary of each.
\end{theorem}

\begin{proof}
We only record the idea from the course.  One
uses degree or mod-two intersection number
with paths crossing the hypersurface.  The
intersection parity changes when crossing
$M$, and this produces two complementary
regions.  A local chart shows that there are
at least two local sides, and the global
intersection argument shows there are no more
than two.
\end{proof}

\section{Borsuk--Ulam theorem}
\label{section-borsuk-ulam}

\begin{theorem}[Borsuk--Ulam]
\label{theorem-borsuk-ulam}
Every continuous map $f:S^n\to\mathbb R^n$
identifies a pair of antipodal points: there
exists $x\in S^n$ such that $f(x)=f(-x)$.
\end{theorem}

\begin{proof}
The smooth case follows from degree theory by
considering the odd map
$g(x)=f(x)-f(-x)$.  If $g$ never vanished,
then $g/\|g\|:S^n\to S^{n-1}$ would give an
odd map to a lower-dimensional sphere, which
contradicts the corresponding degree parity
statement.  The continuous case follows by
approximation.
\end{proof}

\section{Poincare--Hopf theorem}
\label{section-Poincare-Hopf-theorem}

\begin{definition}
\label{definition-isolated-zero-vector-field}
Let $V\in\mathfrak X(X)$ be a vector field.
A point $x\in X$ with $V(x)=0$ is an
{\it isolated zero} if there is a
neighborhood $U$ of $x$ such that
$V(y)\neq0$ for all $y\in U\setminus\{x\}$.
\end{definition}

\begin{definition}
\label{definition-index-vector-field}
Let $x$ be an isolated zero of $V$ on an
$n$-manifold.  Choose coordinates identifying
a small ball around $x$ with $B_r(0)$.
The {\it index} of $V$ at $x$ is
$$
\operatorname{ind}(V,x)
=
\deg\left(
\frac{\widetilde V}{\|\widetilde V\|}
:\partial B_r(0)\to S^{n-1}
\right).
$$
\end{definition}

\begin{theorem}[Poincare--Hopf]
\label{theorem-Poincare-Hopf}
Let $X$ be a compact manifold without
boundary and let $V\in\mathfrak X(X)$ have
isolated zeroes.  Then
$$
\sum_{x\in Z(V)}\operatorname{ind}(V,x)
=\chi(X).
$$
\end{theorem}

\begin{proof}
We give the course-level summary.  Choose a
triangulation or a Morse function and build a
vector field whose zeroes correspond to cells
or critical points.  The local index at each
zero gives the sign with which the cell
contributes.  The sum of these signed local
contributions is the Euler characteristic.
A homotopy argument shows that the sum is
independent of the chosen vector field.
\end{proof}

\section{Hopf theorem}
\label{section-Hopf-theorem}

\begin{theorem}[Hopf]
\label{theorem-Hopf}
Let $X$ be a compact, oriented, connected
$n$-manifold without boundary.  Let
$f,g:X\to S^n$ be smooth maps.  Then
$$
f\sim g
\quad\Longleftrightarrow\quad
\deg(f)=\deg(g).
$$
\end{theorem}

\begin{proof}
One direction is homotopy invariance of the
degree.  Conversely, choose a regular value
$p\in S^n$ for both maps.  The preimages
$f^{-1}(p)$ and $g^{-1}(p)$ are finite signed
sets.  Equality of degrees says that these
signed sets have the same algebraic count.
Using paths and local modifications, one
cancels pairs of opposite signs and moves the
remaining points so that the two maps have
matching inverse images of $p$.  Then the
maps are homotopic on the complement of small
balls and the local degree data fills in the
balls.
\end{proof}

\section{Pontryagin--Thom construction}
\label{section-pontryagin-thom}

\noindent
The Pontryagin--Thom idea is to convert a
submanifold with framing data into a map to a
sphere.  If $W^k\subset M^n$ has a framed
normal bundle, choose a tubular neighborhood
$\nu(W)\cong W\times\mathbb R^{n-k}$.  Collapse
the complement of this neighborhood to the
base point of $S^{n-k}$, and use the framing
to map each normal disk radially to the
sphere.  This produces a map
$$
M\to S^{n-k}.
$$
Cobordism of framed submanifolds corresponds
to homotopy of the resulting maps.

\section{Morse functions}
\label{section-morse-functions}

\begin{definition}
\label{definition-morse-function}
Let $M$ be a smooth manifold.  A smooth
function $f:M\to\mathbb R$ is a
{\it Morse function} if every critical point
is nondegenerate, i.e. the Hessian at that
point is a nondegenerate quadratic form.
\end{definition}

\begin{theorem}[Morse lemma]
\label{theorem-morse-lemma}
Let $p$ be a nondegenerate critical point of
$f:M\to\mathbb R$.  Then there are local
coordinates around $p$ in which
$$
f=f(p)-x_1^2-\cdots-x_\lambda^2
+x_{\lambda+1}^2+\cdots+x_n^2.
$$
The number $\lambda$ is the {\it Morse index}
of $p$.
\end{theorem}

\begin{proof}
The proof is a smooth diagonalization of the
quadratic part of $f$ near the critical
point.  After subtracting the constant
$f(p)$, the first-order terms vanish.  The
nondegenerate Hessian can be diagonalized,
and the remaining higher-order terms are
absorbed by a smooth change of coordinates.
\end{proof}

\begin{theorem}
\label{theorem-morse-functions-generic}
Morse functions are generic in
$C^\infty(M,\mathbb R)$.
\end{theorem}

\begin{proof}
A function fails to be Morse when its 1-jet
meets the zero section in a nontransverse way,
or equivalently when the corresponding
2-jet lies in the degenerate Hessian locus.
This is a jet transversality condition, so
Thom transversality gives the result.
\end{proof}

\section{Handle decomposition}
\label{section-handle-decomposition}

\noindent
Morse theory describes the topology of a
manifold by studying sublevel sets
$$
M^a=\{x\in M:f(x)\leq a\}.
$$
If there are no critical values in
$[a,b]$, then $M^a$ and $M^b$ are
diffeomorphic.  If one crosses a single
critical point of index $\lambda$, then the
sublevel set changes by attaching a
$\lambda$-handle
$$
D^\lambda\times D^{n-\lambda}.
$$

\begin{theorem}
\label{theorem-handle-attachment}
Let $f:M\to\mathbb R$ be Morse and suppose
$c$ is a critical value with one critical
point of index $\lambda$.  For small
$\varepsilon>0$, the sublevel set
$M^{c+\varepsilon}$ is obtained from
$M^{c-\varepsilon}$ by attaching a
$\lambda$-handle.
\end{theorem}

\begin{proof}
This is the local normal form from the Morse
lemma together with the gradient flow.  The
negative directions give the core
$D^\lambda$, and the positive directions give
the cocore $D^{n-\lambda}$.
\end{proof}

\section{Cancellation and four-manifolds}
\label{section-cancellation-four-manifolds}

\noindent
Handle decompositions are not unique.  A
central operation is cancellation: under the
right incidence condition, a
$\lambda$-handle and a $(\lambda+1)$-handle
can cancel each other.  Geometrically, this
corresponds to modifying the Morse function
so that a pair of critical points disappears.

\begin{theorem}[Cancellation principle]
\label{theorem-cancellation-principle}
If the attaching sphere of a
$(\lambda+1)$-handle intersects the belt
sphere of a $\lambda$-handle transversely in
one point, then the two handles cancel.
\end{theorem}

\begin{proof}
We only record the geometric idea.  The
single transverse intersection means that the
gradient trajectory between the two critical
points is unique and has the correct
orientation.  One can then change the Morse
function in a neighborhood of this trajectory
and remove the two critical points.
\end{proof}

\noindent
The course ended by pointing toward the
classification of 4-manifolds.  In dimension
4, handle decompositions are especially rich:
0-, 1-, 2-, 3-, and 4-handles interact with
intersection forms, framings, and Kirby
calculus.  This is the bridge from Morse
theory and differential topology to the
classification problems for smooth
4-manifolds.

\bibliography{my}
\bibliographystyle{amsalpha}

\end{document}
